%% kapa: exam=1 round=26 year=2015 period=2 form=A subject=appraisal_law subject_ko=부동산관계법규 seq=25 number=25
%% source: Q-Net 감정평가사 기출문제 (cst00309) artlSeq=5138001 file=extracted/1008/5138001/제26회_감정평가사 1차_기출/제26회_감정평가사 1차_기출/2교시_감평_A형.pdf
%% conversion: kordoc 4.24.51 -> tools/gen1.py (auto); review=proofread-sonnet; flags=none
\begin{question}{25}
국유재산법령상 행정재산에 관한 설명으로 옳은 것은?
\begin{choices}
  \item 중앙관서의 장은 행정재산으로 사용하기로 결정한 날부터 5년이 지난 날까지 행정재산으로 사용되지 아니한 경우 지체 없이 그 용도를 폐지하여야 한다.
  \item 행정재산은 처분하지 못하므로, 사유재산과 교환하여 그 교환받은 재산을 행정재산으로 관리하려는 경우에도 교환할 수 없다.
  \item 행정재산의 관리수탁자는 위탁받은 재산의 연간 관리현황을 감사원에 보고하여야 한다.
  \item 행정재산을 사용허가하려는 경우 수의(隨意)의 방법을 우선적으로 사용하여야 한다.
  \item 중앙관서의 장은 행정재산에 대하여 일반경쟁입찰을 두 번 실시하여도 낙찰자가 없는 재산에 대하여는 세 번째 입찰부터 최초 사용료 예정가격의 100분의 10을 최저한도로 하여 매회 100분의 5의 금액만큼 그 예정가격을 낮추는 방법으로 조정할 수 있다.
\end{choices}
\end{question}
\answer{1}
